\begin{textsample}{POS Dim 3 – ms – Score 19.00 – ms/ms\_inf\_e\_ef\_28.txt}  \label{ex:f3_pos_232}
7.4.1 O Eu, o Outro e o Nós A interação dentro da instituição de Educação \textbf{Infantil} possibilita à \textbf{criança} o desenvolvimento das linguagens, as diferentes maneiras de se expressar, o estabelecimento de vínculos, a solidariedade, a aprendizagem e a autonomia.
A construção da identidade pela \textbf{criança} ocorre por meio das relações sociais que são estabelecidas no convívio com a família, com outras \textbf{crianças} e \textbf{adultos}.
Nesse sentido, as DCNEI ( BRASIL, 2009 ) preconizam que a \textbf{criança} é : > sujeito histórico e de direitos que, nas interações, relações e práticas cotidianas que vivencia, constrói sua identidade pessoal e coletiva, \textbf{brinca}, imagina, fantasia, deseja, aprende, observa, experimenta, narra, questiona e constrói sentidos sobre a natureza e sociedade produzindo cultura.
Desse modo, é a partir das experiências vivenciadas que a \textbf{criança} compreende e modifica o mundo e a si mesma.
A instituição é um espaço importante para as interações com o outro, sendo fundamental a intencionalidade educativa nos contextos de aprendizagem, para que os profissionais permitam os diálogos entre as \textbf{crianças}, o interesse, a \textbf{curiosidade}, os questionamentos, os \textbf{cuidados} com seu corpo e ambiente, a exploração dos espaços e materiais, a construção da autonomia e de aprendizagens.
Assim, esse campo propõe um trabalho a partir das situações do cotidiano que perpassam as relações humanas, integrando -o a outros campos.
Dessa forma : > Cada campo de experiência oferece um conjunto de objetos, situações, imagem e linguagens, referidos aos sistemas simbólicos de nossa cultura, capazes de evocar, estimular, acompanhar aprendizagens \textbf{progressivamente} mais seguras ( BARBOSA ; FARIA, FINCO, 2015, p. 54-55 ).
A \textbf{criança} desenvolve autonomia, autocuidado e identidade na interação com seus pares e na participação de atividades significativas.
Dessa forma, a escola é um ambiente com a capacidade de \textbf{acolher} as diferenças e de promover as potencialidades das \textbf{crianças}, contribuindo, assim, para suas aprendizagens e seu desenvolvimento. 7.4.2 Corpo, \textbf{Gestos} e Movimentos A infância é um período no qual os movimentos corporais ocupam grande parte do tempo das \textbf{crianças}.
Desde que nascem, os \textbf{bebês} fazem uso do corpo para conhecerem a si mesmos e o mundo ; estabelecendo relações com \textbf{adultos} e outras \textbf{crianças} tomam consciência do seu corpo e das possibilidades de movimento e sensações que ele proporciona.
Assim, eles utilizam o corpo como elemento central de comunicação e expressão para com o outro e com o mundo.
Nessa perspectiva : > A ação do corpo propicia viver emoções e sensações prazerosas, de relaxamento e de tensão, mas também a satisfação do controle dos \textbf{gestos}, na coordenação com os outros ; consente experimentar potencialidades e limites do próprio aspecto físico, desenvolvendo ao mesmo tempo, a consciência dos riscos dos movimentos incontrolados ( BARBOSA ; FARIA, FINCO, 2015, p. 58 ).
Nesse campo, as experiências motoras nas \textbf{brincadeiras}, danças e dramatizações possibilitam às \textbf{crianças} identificarem possibilidades na exploração de \textbf{gestos} e movimentos corporais.
Elas descobrem, \textbf{brincam}, manifestam sentimentos, sensações, emoções, \textbf{desejos} e medos, experiências que vivenciam sozinhas ou na interação com os outros.
Enfim, o corpo “ fala ” e para que a \textbf{criança} vivencie, interaja e aprenda, utilizando -o como apoio, é necessário que seja estimulada a ampliar seus movimentos e o seu conhecimento de mundo, com ações planejadas e a organização de situações que promovam o seu desenvolvimento. 7.4.3 Traços, \textbf{Sons}, Cores e Formas As orientações descritas nas DCNEI ( 2009 ) apontam para a elaboração de propostas pedagógicas nas quais os princípios estéticos e as diferentes manifestações artísticas e culturais devem ser considerados.
O artigo 9º dispõe que a prática pedagógica deve assegurar experiências que promovam o relacionamento e a interação das \textbf{crianças} com diversificadas manifestações de música, artes \textbf{plásticas} e gráficas, cinema, fotografia, dança, teatro, poesia e literatura.
A imersão dessas manifestações artísticas no currículo da Educação \textbf{Infantil} cria possibilidades para as \textbf{crianças} investigarem o mundo que está a sua \textbf{volta} e a si mesmas de forma poética e estética, desenvolvendo, assim, a sensibilidade e a expressividade diante da vida, pois, segundo Gobbi ( 2010, p. 3 ), “ a dimensão \textbf{lúdica} e a dimensão estética são condições fundamentais para a formação humana ”.
Trabalhar com o campo “ Traços, \textbf{sons}, cores e formas ” significa garantir à \textbf{criança} experiências artísticas significativas e inseri -la em um universo que vem se constituindo ao longo da existência humana.
É possibilitar a apropriação e a expressão por meio de linguagens diversificadas, cada qual com saber consolidado, com especificidades próprias e que, além dessas especificidades, cumprem o papel de “ contribuir para transformar o arranjo da consciência dos homens, conferindo -lhes novas formas de apreensão do real ” ( FERREIRA, 2010, p. 14 ), bem como da transformação da realidade.
Assim, na Educação \textbf{Infantil} : > O mundo é multicolorido, multiforme e diverso em suas \textbf{texturas}.
Incentivar que as \textbf{crianças} fiquem mais atentas para as cores, formas e \textbf{texturas} do mundo – tanto da natureza quanto aos bens culturais – é um grande estímulo para elas e uma oportunidade de sensibilizar seu olhar e ampliar o seu repertório.
Não importa se a \textbf{criança} está usando formas específicas e cores tais e quais – menos ainda se sabem seus nomes corretamente.
Nos \textbf{interessa}, sim, que suas expressões ( desenhos, \textbf{pinturas}, colagens ) possam explorar diferentes formas, cores e \textbf{texturas}, ampliando seu acervo de sensações e de imagens, combinando -as entre si, de maneira própria, autoral ( BRASIL, 2006, p. 67 ).
Dessa forma, quanto mais amplo for o repertório da \textbf{criança} em relação aos bens culturais, mais possibilidades ela terá para aprender a se expressar de forma autoral pelo desenho, \textbf{pintura}, modelagem, colagem, etc.
Quanto mais ricas forem as experiências com a diversidade de linguagens, maiores serão as possibilidades de expressão e criação.

% matched lemmas: acolher, adulto, bebê, brincadeira, brincar, criança, cuidado, curiosidade, desejo, gesto, infantil, interessar, lúdico, pintura, plástico, progressivamente, som, textura, volta
\end{textsample}
